\begin{table}[H]
\centering\small\TableSetup
\caption{In den Entwicklungsstudien verwendete historische arithmetische Ersatzgrößen. Bei Zählwerten kommen MAC-äquivalente und skalare Ersatzgrößen zum Einsatz. Tabelle~\ref{tab:cost:operation_extensions} enthält die zusätzlichen Regeln, die bei der aktuellen Bewertung angewendet werden.}
\label{tab:flop_formulas_landscape}
\begin{tabularx}{\linewidth}{@{}L{3.5cm}Y@{}}
\toprule
Operatorfamilie & Zählweise und Konvention \\
\midrule
Conv, ConvInteger, QLinearConv & $NH_{out}W_{out}C_{out}(K_hK_wC_{in}/G)$. Bei explizitem Bias wird $N_{out}$ addiert. Eine MAC-äquivalente Einheit pro Kernel-Produkt pro Ausgangskanal. \\
ConvTranspose & Gleiche Ausdrucks- und Bias-Konvention wie bei Conv. Historisches, auf dichten Ausgabewerten basierendes Surrogat, keine operatorgenaue MAC-Zählung. \\
MatMul, MatMulInteger, QLinearMatMul & $BMNK$. Formen werden durch Inferenz aus Werten oder Initialisierern abgeleitet. Pro Produkt eine MAC-äquivalente Einheit. \\
Gemm & $BMNK$. Bei explizitem Bias wird $N_{out}$ addiert. \\
Add, Sub, Mul, Div & $N_{out}$. Eine skalare Proxy-Einheit pro Ausgabekomponente. \\
Batch-Normalisierung & $4N_{out}$. Fester Standardwert pro Ausgabeelement. \\
ReLU, LeakyReLU, Sigmoid, tanh, Softsign, Softplus, GELU, Clip & $N_{out}$. Fester Standardwert von einer Proxy-Einheit pro Ausgabeelement. \\
Resize, Upsample, Interpolate & $c_{mode}N_{out}$ mit $c_{mode}=1,8,16$ für den Nächste-Nachbar-, linearen und kubischen Modus. Unbekannte Modi werden standardmäßig auf „eins“ gesetzt. \\
Sonstige Operatoren & Null, sofern nicht in aufgeführte Primitive zerlegt. Dazu gehören „Fused Attention“, „Pooling“, „Softmax“, „Layer Normalisation“, „Transpose“ und „Reshape“. Ihre Auswirkungen auf Laufzeit und Speicherbedarf bleiben beobachtbar. \\
\bottomrule
\end{tabularx}
\end{table}
